\section{A.5. Four points, Klein groups, and the octahedron}
We work first over an algebraically closed field $k$ of characteristic different from $2$. In cote 67/084--089, $V$ denotes a Klein four group. We write $V^\circ=V\setminus\{1\}$ and
\[
 S(V)=\bigcup_{v\in V^\circ}(\mathbf P^1)^v.
\]
The word read tentatively as ``r\'eelle'' in 67/086 is incompatible with its stated algebraically closed base; none of the arguments below uses that word. The object is a projective line over $k$.

\paragraph{Why characteristic $2$ is excluded.}
In characteristic $2$, $i^2=-1$ implies $(i-1)^2=0$, hence $i=1$ in a field. The displayed six-set then collapses to $\{0,1,\infty\}$, and $\{z,-z,1/z,-1/z\}$ has only two elements. This is not merely a defect of that coordinate choice: over an algebraically closed field of characteristic $2$, the four translations $z\mapsto z+a$ for $a\in\mathbf F_4$ give a faithful Klein action, but each nonidentity translation has the single geometric fixed point $\infty$. Its union of fixed points therefore has cardinality one, not six. The octahedral fixed-pair argument cannot extend to this characteristic. These identities give an exact algebraic check of the excluded edge case, without relying on an odd-prime enumeration. The accompanying exact certificate \texttt{code/part\_a567\_char2\_edge.py} checks all sixteen products of the four translations and all five rational projective points over $\mathbf F_4$; its output is attached.

\begin{proposition}[Proved: the normalizer, including characteristic $3$ and $5$]\label{A5-normalizer}
Every faithful action of $V$ on $\mathbf P^1_k$ is conjugate to
\[
 V_0=\{z,\ -z,\ 1/z,\ -1/z\}.
\]
Its centralizer in $\operatorname{PGL}_2(k)$ is $V_0$, and its normalizer is isomorphic to $\mathfrak S_4$. The maps
\[
 V\longmapsto N(V),\qquad \Gamma\longmapsto
 \text{the unique normal Klein subgroup of }\Gamma
\]
are inverse bijections between Klein subgroups and subgroups isomorphic to $\mathfrak S_4$, in every characteristic different from $2$.
\end{proposition}
\begin{proof}
An involution has two distinct fixed points and is conjugate to $z\mapsto-z$. A different commuting involution must interchange those fixed points: the transformations preserving them individually are $z\mapsto az$, whose only nonidentity involution is $-z$. It therefore has form $z\mapsto a/z$, and rescaling $z$ makes it $1/z$. Here we use existence of a square root of $a$.

A transformation commuting with $-z$ is $az$ or $a/z$; commuting also with $1/z$ forces $a^2=1$. Thus the centralizer is $V_0$. Conjugation embeds $N(V_0)/V_0$ into $\operatorname{Aut}(V_0)\simeq\mathfrak S_3$. The transformations
\[
 h(z)=iz,\qquad b(z)=\frac{z+1}{z-1},\qquad i^2=-1,
\]
normalize $V_0$ and induce two distinct transpositions of its three nonidentity elements. Consequently $|N(V_0)|=24$.

To identify this group without a characteristic-dependent classification of finite projective groups, label the six fixed points $0,\infty,1,-1,i,-i$ and pair them as
\[
 \{0,\infty\},\quad\{1,-1\},\quad\{i,-i\}.
\]
The four pairs of complementary triples choosing one point from each pair are the four antipodal pairs of faces. On this four-element set the permutations induced by $h$ and $b$ generate $\mathfrak S_4$: one may label the triples by $(\epsilon_1,\epsilon_2,\epsilon_3)\in\{\pm1\}^3$ modulo simultaneous negation; $V_0$ acts regularly, and $h,b$ induce the full permutation group on its three nontrivial translations. Equivalently, the resulting affine group is $\mathbf F_2^2\rtimes\operatorname{GL}_2(\mathbf F_2)$ of order $24$. For completeness, a permutation preserving the three antipodal pairs acts on the face labels as a permutation of the three coordinates followed by a sign change. If it acts trivially on $\mathbf F_2^3/\langle(1,1,1)\rangle$, its coordinate permutation is trivial (the three coordinate vectors remain distinct in this quotient), and its sign change is either trivial or simultaneous negation. The latter is the antipodism of the six vertices. That antipodism is not projective: a projective map interchanging $0,\infty$ and $1,-1$ must be $-1/z$, which fixes $i$ instead of sending it to $-i$. Hence the action of $N(V_0)$ on the four face pairs is faithful, and its order proves $N(V_0)\simeq\mathfrak S_4$.

Every abstract $\mathfrak S_4$ has its unique normal Klein subgroup (the identity and the three double transpositions). If $\Gamma\subset\operatorname{PGL}_2(k)$ is such a group and $V$ is that subgroup, then $\Gamma\subset N(V)$ and both groups have order $24$, so they are equal. This proves both inverse assertions. All six displayed points remain distinct in characteristics $3$ and $5$; no step divided by $3$ or $5$.
\end{proof}

\begin{theorem}[Refuted in characteristic $5$; corrected proposition of 67/086]\label{A5-six-points}
If $\operatorname{char}k\ne2,5$, a six-point set projectively equivalent to
$S_0=\{0,\infty,\pm1,\pm i\}$ determines a unique Klein subgroup $V$ with $S(V)=S_0$, and its full projective stabilizer is $N(V)\simeq\mathfrak S_4$.
In characteristic $5$, the same set is $\mathbf P^1(\mathbf F_5)$, its stabilizer has order $120$, and exactly five Klein subgroups have this fixed-point set. Retaining the antipodal pairing determined by $V$ restores the correspondence in characteristic $5$.
\end{theorem}
\begin{proof}
Call an unordered quartet harmonic if one, hence every, ordering has cross-ratio in $\{-1,2,1/2\}$. For each antipodal pair of $S_0$, its complementary quartet is harmonic. All nonantipodal pairs are equivalent under the normalizer from Proposition~\ref{A5-normalizer}: the corresponding assertion is transitivity on the twelve edges of the octahedron, or follows directly from the two displayed permutations. For the pair $\{0,1\}$, the complementary quartet $\{\infty,-1,i,-i\}$ has a cross-ratio equal to $i$ or $-i$, depending on the ordering. Such a value is harmonic only if $i$ or $-i$ is $-1,2$, or $1/2$. Squaring shows respectively $2=0$, $5=0$, or $5/4=0$. Thus outside characteristics $2,5$, precisely the three antipodal pairs have harmonic complements. In characteristic $3$ the three harmonic values coincide, but the preceding distinction still holds.

Any projective stabilizer of $S_0$ must therefore preserve its three pairs. For a chosen pair there is a unique nonidentity projective involution fixing its two points and interchanging the points of each other pair: after sending that pair to $0,\infty$, it is $z\mapsto-z$. The three involutions reconstruct $V_0$. Hence the stabilizer normalizes $V_0$, proving that it is $N(V_0)$ and proving uniqueness of $V$.

In characteristic $5$ choose $i=2$. Then $S_0=\mathbf P^1(\mathbf F_5)$. Every projective map stabilizing this set is defined over $\mathbf F_5$, because its images on the three points $\infty,0,1$ determine it uniquely and those images are rational. Conversely every element of $\operatorname{PGL}_2(\mathbf F_5)$ preserves the set. The order is
\[
 |\operatorname{PGL}_2(\mathbf F_5)|=5(5^2-1)=120.
\]
All Klein subgroups over $k$ are conjugate by Proposition~\ref{A5-normalizer}. If two have the same fixed-point set, any conjugating projective map stabilizes that set. Hence the fibre over $S_0$ is one orbit of its stabilizer, with isotropy $N(V_0)$ of order $24$, and has $120/24=5$ elements. Finally, the reconstruction from the three distinguished pairs works in characteristic $5$ as well, so retaining the pairing repairs the statement.

The exact program \texttt{code/part\_a567\_checks.py} enumerates all $6\cdot5\cdot4=120$ possible images of the projective triple for each odd prime below $100$. It checks the full stabilizer and every Klein subgroup whose fixed-point union is the six-set; in characteristic $5$ the outputs are respectively $120$ and $5$. This finite certificate supplements, and does not replace, the all-characteristic proof.
\end{proof}

\begin{proposition}[Proved: the face stabilizer correction on 67/085]\label{A5-faces}
A single face of the oriented octahedron has stabilizer of order $3$ in $N(V)$. An unordered pair of opposite faces has stabilizer $\mathfrak S_3$, complementary to $V$. Thus the splitting on 67/085 belongs to a pair of opposite faces.
\end{proposition}
\begin{proof}
The group of order $24$ acts transitively on the eight faces, so a single stabilizer has order $3$. It acts as the full symmetric group on the four pairs of opposite faces, so a pair stabilizer is $\mathfrak S_3$. The Klein subgroup acts regularly on those four pairs, and thus meets such a stabilizer trivially. Their product has order $4\cdot6=24$, which is the required splitting.
\end{proof}

\paragraph{The two orientations on 67/085--086.}
The six labelled vertices define an abstract octahedron, also when $k$ has positive characteristic. The two normalizer generators above act as rotations: $h$ is a quarter-turn about the $0,\infty$ axis, while $b$ exchanges $0,-1$, exchanges $\infty,1$, and exchanges $i,-i$, hence is an edge-axis half-turn. Thus every change of standard coordinate preserves the orientation in which $(0,1,i)$ is a positive face. Replacing $i$ by $-i$ reverses that orientation. This proves the claimed canonical identification of the two orientations with the two roots of $-1$; no ordering or real structure on $k$ is used.

\begin{theorem}[Proved: the groupoid theorem of 67/089]\label{A5-groupoid}
Let $k$ be algebraically closed, $\operatorname{char}k\ne2$, and fix a choice of $i^2=-1$. Let $\mathcal C_k$ be the groupoid whose objects are projective lines $X/k$ with an embedded four-element set $I$, and whose arrows are projective isomorphisms carrying the embedded set bijectively to the embedded set. For $I$, let $J(I)$ be its three partitions into two pairs and let $Z(J(I))$ be the projective line with those three labels as marked points. Then $\mathcal C_k$ is equivalent to the groupoid of pairs
\[
 (I,\zeta),\qquad \zeta\in Z(J(I))\setminus J(I),
\]
with arrows the bijections of $I$ whose induced maps of three-pointed lines carry $\zeta$ to $\zeta'$. The torsor-twist construction on 67/088--089 is an inverse. No exception at characteristic $3$ or $5$ is needed.
\end{theorem}
\begin{proof}
The double transpositions of $I$ extend uniquely to projective automorphisms of $X$. For example, with $I=\{\infty,0,1,\lambda\}$ they are obtained by solving the prescribed images of three points, and the fourth image follows from invariance of cross-ratio under double transpositions. These automorphisms form a faithful Klein group $V(I)$. Its action on $I$ is free and transitive.

The quotient $X/V(I)$ is a projective line with three branch points indexed by $V(I)^\circ$, equivalently by $J(I)$. For the standard Klein action its quotient function is
\[
 \lambda=\frac{(z^2+1)^2}{4z^2},
\]
whose three branch values are $0,1,\infty$, with respective fixed-point pairs $\{\pm i\},\{\pm1\},\{0,\infty\}$. Its degree is four, and the group action proves it is the quotient, including in characteristics $3$ and $5$. The orbit $I$ maps to a point $\zeta$ away from these branch values. Since a map of projective lines is determined by the images of three distinct points, the three-pointed quotient is canonically $Z(J(I))$.

For clarity we construct the inverse in the manuscript's notation. Fix a standard four-element set $I_0$ of antipodal face pairs for $V_0$, and use Proposition~\ref{A5-normalizer} to identify its projective normalizer with $\mathfrak S_{I_0}$. Transporting the standard projective line along the torsor $\operatorname{Isom}(I_0,I)$ gives $X(I)$ with its $V(I)$ action and its canonical quotient $Z(J(I))$. The fibre
\[
 P=X(I)_\zeta
\]
is a $V(I)$-torsor. Let $T=\operatorname{Isom}_{V(I)}(P,I)$, also a $V(I)$-torsor. Form the contracted product $X=X(I)\times^{V(I)}T$. Evaluation identifies its fibre over $\zeta$ with $I$: the relation $(p,t)\sim(vp,t\circ v^{-1})$ leaves $t(p)$ unchanged. This gives the required embedded four-set.

Conversely, for an original $(X,I)$, any $V(I)$-equivariant identification of the standard Klein action with its action on $X$ identifies the quotients, since its ambiguity is the centralizer $V(I)$ and that acts trivially on the quotient. Its restriction to the free fibre gives exactly the torsor $T$ above. Contracted-product descent then recovers $X$ and the embedding of $I$. This argument also identifies every arrow: changing a chosen standard identification by $V(I)$ changes the fibre trivialization by precisely the contracted-product relation. Thus the two constructions are quasi-inverse on objects and on morphisms.
\end{proof}

\begin{proposition}[Proved with explicit distinction: L7 as a stack and as a coarse space]\label{A5-moduli}
Set $R=\Z[1/2,i]$ and $U=\operatorname{Spec}R[\lambda,\lambda^{-1},(\lambda-1)^{-1}]$. The family version of Theorem~\ref{A5-groupoid} gives
\[
 \mathcal M_{0,[4],R}\simeq[U/\mathfrak S_4],
\]
where the action factors through $\mathfrak S_4/V\simeq\mathfrak S_3$ and the usual six anharmonic transformations of $\lambda$. Its coarse moduli space is $\mathbf A^1_R$, with coordinate
\[
 J=\frac{(\lambda^2-\lambda+1)^3}{\lambda^2(\lambda-1)^2}.
\]
In particular the groupoid theorem does not identify the unordered moduli stack with its coarse space: generic inertia is $V$.
\end{proposition}
\begin{proof}
An ordered four-pointed genus-zero curve has a unique coordinate sending its first three sections to $\infty,0,1$; the fourth is $\lambda$. This proves that the ordered moduli functor is $U$. Relabelling gives the stated $\mathfrak S_4$ action, whose kernel is the double-transposition Klein subgroup. Forgetting the ordering is descent along the finite \'etale torsor of orderings, hence gives the quotient stack. The explicit normalizer and torsor constructions above are defined over $R$ and commute with this descent. This supplies the family argument that an equivalence on algebraically closed points alone would not supply.

Let $A=R[\lambda,\lambda^{-1},(\lambda-1)^{-1}]$. The function $J$ is invariant under $\lambda\mapsto1-\lambda$ and $\lambda\mapsto1/\lambda$, so under the full group. The equation
\[
 (\lambda^2-\lambda+1)^3-J\lambda^2(\lambda-1)^2=0
\]
is monic in $\lambda$. Its transformed versions show that $\lambda^{-1}$ and $(\lambda-1)^{-1}$ are integral over $R[J]$, and thus $A$ is finite over $R[J]$. Over the fraction field of $R$, the rational function $J$ has degree six. The six anharmonic transformations are distinct and fix it, so its field is exactly the invariant field. Therefore $A^{\mathfrak S_3}$ is integral over $R[J]$ and lies in its fraction field. Since $R[J]$ is integrally closed, $A^{\mathfrak S_3}=R[J]$. The finite-group quotient is consequently $\mathbf A^1_R$. The same degree-six argument works separately over algebraically closed residue fields of characteristics $3$ and $5$: the numerator and denominator remain coprime and the six rational transformations remain distinct. Finally, the kernel $V$ fixes every point of $U$, giving the asserted generic inertia.
\end{proof}


\paragraph{The base in L7.}
The quotient-stack presentation and its coarse invariant already descend
to $\Z[1/2]$. Indeed ordered four-pointed curves, their relabelling
action, and the displayed function $J$ are defined over that ring; the
same integrality and normality proof computes its invariant ring.
Adjoining $i$ is needed for the particular displayed octahedral
matrices, not for the existence of this moduli stack or its coarse
space. No averaging by $|\mathfrak S_3|$ is used at characteristic $3$.

\paragraph{Exceptional inertia and the characteristic-$3$ qualification.}
When the characteristic is different from $2,3$, the harmonic values
$\{-1,2,1/2\}$ have stabilizer of order $2$ in the anharmonic
$\mathfrak S_3$, and the roots of $\lambda^2-\lambda+1$ have
stabilizer of order $3$. Their inverse images in $\mathfrak S_4$
are respectively a dihedral group of order $8$ and $\mathfrak A_4$;
the other inertia groups are $V$. In characteristic $3$ these
exceptional loci coalesce at $\lambda=-1$, which all six
anharmonic transformations fix. Its inertia is therefore
$\mathfrak S_4$ of order $24$, while all other inertia groups are
$V$. This follows directly by solving the fixed-point equations
of $1-\lambda$, $1/\lambda$, and $1/(1-\lambda)$. Thus the
earlier input event giving exceptional orders $8$ and $12$ needs
this characteristic-$3$ qualification, although its groupoid
equivalence is valid there.
